\documentclass[12pt]{article}
\usepackage{amsmath, amssymb}
\usepackage{geometry}
\geometry{margin=1in}

\title{Why $\bar{y}_1 - \bar{y}_0$ is Unbiased for the ATE under Random Assignment}
\author{}
\date{}

\begin{document}
\maketitle

\textbf{Claim.} Under random assignment, $E[\bar{y}_1 - \bar{y}_0] = ATE$.

\bigskip

\textbf{Definition of unbiased.} An estimator $\hat{\theta}$ is \textbf{unbiased}
for parameter $\theta$ if:
\[
  E[\hat{\theta}] = \theta
\]
Here $\hat{\theta} = \bar{y}_1 - \bar{y}_0$ and $\theta = ATE = E[y(1) - y(0)]$.

\bigskip

\textbf{Step 1: What do we observe in each group?}

\medskip
Treated individuals reveal $y(1)$; control individuals reveal $y(0)$:
\[
  E[\bar{y}_1] = E[y \mid x=1] = E[y(1) \mid x=1]
\]
\[
  E[\bar{y}_0] = E[y \mid x=0] = E[y(0) \mid x=0]
\]

\bigskip

\textbf{Step 2: Apply random assignment.}

\medskip
Random assignment means $x \perp \bigl(y(0),\, y(1)\bigr)$: knowing whether
someone was treated tells you nothing about their potential outcomes. Therefore:
\[
  E[y(1) \mid x=1] = E[y(1)], \qquad E[y(0) \mid x=0] = E[y(0)]
\]

\bigskip

\textbf{Step 3: Substitute and conclude.}

\begin{align*}
  E[\bar{y}_1 - \bar{y}_0]
    &= E[y(1) \mid x=1] - E[y(0) \mid x=0] & \text{(Step 1)}\\
    &= E[y(1)] - E[y(0)]                    & \text{(Step 2)}\\
    &= E[\,y(1) - y(0)\,]                   & \text{(linearity of expectation)}\\
    &= ATE                                   & \checkmark
\end{align*}

\bigskip

\textbf{Intuition.} The fundamental problem of causal inference is that we never
observe both $y(1)$ and $y(0)$ for the same person. But with
$x \perp (y(0), y(1))$:
\begin{itemize}
  \item The treated group is a \emph{random sample} of the population
        $\Rightarrow$ their average $y(1)$ represents \textbf{everyone's} $y(1)$.
  \item The control group is a \emph{random sample} of the population
        $\Rightarrow$ their average $y(0)$ represents \textbf{everyone's} $y(0)$.
\end{itemize}
Each group serves as a valid counterfactual for the other, and selection bias
cancels exactly.

\end{document}
